\documentclass[a4paper,12pt]{article}
\usepackage[utf8]{inputenc}
\usepackage[T1]{fontenc}
\usepackage[ngerman]{babel}
\usepackage{amsmath}
\usepackage{amssymb}
\usepackage{enumitem}
\usepackage{geometry}
\geometry{a4paper, margin=2.5cm}
\usepackage{parskip}

\title{Einsendeaufgaben -- Aufgabe 5.2\\[0.5em]
\large Modul 61211: Analysis}
\author{}
\date{}

\begin{document}
\maketitle

\section*{Aufgabe 5.2}

Wir beobachten zunächst, dass $\int_0^\infty g'(x)\,dx$ integrierbar ist, da
\[
  \lim_{a\to\infty} \int_0^a g'(x)\,dx = \lim_{a\to\infty} g(a) - g(0) = -g(0).
\]
Demnach ist auch $\int_0^\infty M|g'(x)|\,dx$ integrierbar.

Wegen $|F(x)| \le M$ folgt $|F(x)g'(x)| \le M|g'(x)|$ für alle $x \in [0,\infty)$. Wir können also das Majorantenkriterium anwenden und schließen, dass $\int_0^\infty F(x)g'(x)\,dx$ integrierbar ist.

Nach der partiellen Integration folgt für $a\to\infty$:
\[
  \int_0^a f(x)g(x)\,dx = F(a)g(a) - F(0)g(0) - \int_0^a F(x)g'(x)\,dx \longrightarrow c \quad \text{mit } c \in \mathbb{R}.
\]
Denn wegen der Beschränktheit von $F$ folgt $F(a)g(a) \to 0$ für $a\to\infty$ und $Fg' \in \mathcal{R}([0,\infty])$.

\end{document}
